\icode{E-EMM-SI-L2-260625-001}
如图所示，垂直纸面方向有两条通电长直导线\(a\)、\(b\)，导线中通有大小相等，方向相反的电流，它们之间的距离为\(L\)。已知通电长直导线周围某点的磁感应强度为\(B = k \frac{I}{r}\)， 其中\(r\)为该点到导线的距离，\(I\)为电流大小。\(O\)为\(a\)、\(b\)连线的中点，\(P\)、\(Q\)为\(a\)、\(b\)连线中垂线上两点， \(P\)、\(Q\)两点到\(a\)、\(b\)的距离均为\(L\)。 已知直导线\(a\)在\(O\)点产生的磁感应强度为\(B_0\)。下列说法正确的是\paren[]

\begin{tikzpicture}
    \draw[dashed] (1.73,0)--(-1.73,0) (0,-1)--(0,1);
    \filldraw[fill=white] (0,1)node{\(\times\)} circle(.15);
    \filldraw[fill=white] (0,-1)node{\(\cdot\)} circle(.15);
    \fill (-1.73,0)node[above,font=\scriptsize]{\(P\)} circle(.05);
    \fill (1.73,0)node[below,font=\scriptsize]{\(Q\)} circle(.05);
    \node at(.4,1){\(a\)};
    \node at(.4,-1){\(b\)};
    \node[font=\scriptsize] at(-.2,-.2){\(O\)};
\end{tikzpicture}
\begin{choices}
    \item 导线\(a\)在\(P\)点产生的磁感应强度大小为\(B_0\)
    \item \(O\)点的磁感应强度大小为0
    \item \(P\)点的磁感应强度大小为\(\frac{B_0}{2}\)
    \item \(P\)、\(Q\)两点的磁感应强度方向相反
\end{choices}
\begin{answer}
    C
\end{answer}